\scanpage{138}
\mspage{V.10}
% archive sheet number: 137
% manuscript lines 1-4
\[
H^i(X_0,\Lambda_{\unsure{\nu|v}})\longrightarrow H^i(X_{\unsure{\bar\eta|\bar y}},\Lambda_\nu)
\]
est un \unsure{isomorphisme} si $i>n+d+1$, un épimorphisme si $i=n+d+1$. L'action du groupe d'inertie \deleted{sur} $I$ sur $H^i(X_{\bar\eta},\Lambda_\nu)$ est triviale si $i\geq n+d+1$.

% manuscript lines 5-10
Procédant par dualité on trouve :

\uline{Corollaire} 2.6. Avec les hypothèses précédentes, l'action de $I$ sur $H^i(X_{\bar\eta},\Lambda_\nu)$ est triviale si $i\notin[n-d,n+d]$.
En particulier, si $d=0$ i.e. si l'ens $Z$ des pts de $X_0$ où $f$ est non lisse est fini, alors l'action de $I$ sur $H^i(X_{\bar\eta},\Lambda_\nu)$ est triviale si $i\ne n$.

% manuscript lines 11-15
\uline{Remarque} 2.7. Lorsque $f:X\mathbin{\unsure{-|\to}}S$ est projectif, on peut trouver une démonstration plus simple de \deleted{\unsure{trivialité}} 2.5 (et \inserted{de} 2.6) \deleted{\unsure{de l'action triviale du groupe d'inertie}}. En effet, \deleted{\unsure{plongeons}} utilisons un plongement $X\subset\mathbb P^r_S$, et soit $Y$ un sous-schéma de $X$, obtenu en coupant $X$ avec $d+1$ \deleted{\unsure{sections}} hyperplans de $\mathbb P^r_S$.
% slanted
\margin{intervertir avec 2.8}
% manuscript lines 16-22
Si ces derniers sont bien choisis, on peut supposer $Y_0$ lisse de codim $d+1$ dans $X_0$, et alors $Y$ est lisse sur $S$, de codim $d$ dans $X$. De plus, le th. de Lefschetz affine SGA 4 XIV implique aisément, \deleted{\unsure{par les}} de façon connue, que l'hom de Gysin
\[
H^i(Y_{\bar\eta},\Lambda_\nu)\mathbin{-}H^{\deleted{\unsure{i}}\inserted{i+2(d+1)}}(X_{\bar\eta},\Lambda_\nu(\deleted{\unsure{+|2|*}}d+1))
\]
